\chapter*{Preface}
\phantomsection
\addcontentsline{toc}{chapter}{Preface}
\markboth{Preface}{Preface}
\bichaptertitle{前言}

\section*{Systematic trading and investing}
\phantomsection
\addcontentsline{toc}{section}{Systematic trading and investing}
\bisectiontitle{系统化交易与投资}

I am very bad at making financial decisions. Like most people I find it difficult to
manage my investments without becoming emotional and behaving irrationally. This is
deeply irritating as I consider myself to be very knowledgeable about finance. I've
voraciously read the academic literature, done my own detailed research, spent 20 years
investing my own money and nearly a decade managing funds for large institutions.

我很不擅长做金融决策。和大多数人一样，我发现自己在管理投资时很难不受情绪影响，也很难始终保持理性。这让我非常沮丧，
因为我一直认为自己对金融相当了解。我大量阅读学术文献，做过细致的个人研究，用自己的钱投资了 20 年，还曾近十年为大型机构管理资金。

So in theory I know what I'm doing. In practice when faced with a decision to buy or sell
a stock things go wrong. Fear and greed wash through my mind, clouding my judgment.
Even if I've spent weeks researching a company it's still hard to click the trade button on
my broker's website. I have to stop myself buying or selling on a whim, based on nothing
more than random newspaper articles or an anonymous blogger's opinion. But then, like
you, I'm only human.

所以，从理论上说，我知道自己在做什么。但在现实中，一旦真正面对买入或卖出股票的决定，事情就会出问题。恐惧与贪婪会涌上心头，
蒙蔽我的判断。即使我已经花了几周时间研究一家公司，在经纪商网站上按下交易按钮时依然很难下决心。我必须不断克制自己，
不要仅仅因为几篇随机的报纸文章，或者某位匿名博主的看法，就一时冲动地买卖。不过说到底，和你一样，我也只是个普通人。

Fortunately there is a solution. The answer is to fully, or partly, systematise your
financial decision making. Creating a trading system removes the emotion and makes it
easier to commit to a consistent strategy.

幸运的是，这个问题有解。答案就是把你的金融决策过程全部或部分系统化。建立一套交易系统可以剥离情绪因素，让你更容易坚持一套一致的策略。

I spent many years managing a large portfolio of trading strategies for a systematic hedge
fund. Unfortunately I didn't have the opportunity to develop and trade systems to look
after my personal portfolio. But after leaving the industry I've been able to make my own
trading process entirely systematic, resulting in significantly better performance.

我曾在一家系统化对冲基金中多年负责管理庞大的交易策略组合。遗憾的是，当时我没有机会为自己的个人投资组合开发并运作交易系统。
但离开这个行业之后，我终于能够把自己的交易流程彻底系统化，而结果是业绩明显改善。

There are many authors and websites offering trading systems. But many of these
``systems'' require subjective interpretation, so they are not actually systematic. Some are
even downright dangerous, trading too quickly and expensively, and in excessive size.
They present you with a single ``one size fits all'' system which won't suit everybody. I
will explain how to develop your own trading system, for your own needs, and which
should not be excessively dangerous or costly to operate.

市面上有很多作者和网站提供各种交易系统。但其中许多所谓的“系统”仍然需要主观解读，因此并不是真正的系统化方法。有些甚至相当危险，
交易过于频繁、成本过高，而且仓位过大。它们通常只提供一种“通用于所有人”的单一系统，但这种系统不可能适合每个人。我会说明如何根据你自己的需求，
开发属于你自己的交易系统，而且它不应当在风险或操作成本上走向极端。

I've found systematic investing to be more profitable, and to require less time and effort.
This book should help you to reap similar benefits.

我发现，系统化投资不仅更有盈利性，而且所需的时间和精力也更少。这本书应当能帮助你获得类似的收益。

\section*{Who should read this book}
\phantomsection
\addcontentsline{toc}{section}{Who should read this book}
\bisectiontitle{谁应该读这本书}

This book is intended for everyone who wishes to systematise their financial decision
making, either completely or to some degree.

本书面向所有希望将自己的金融决策全部或部分系统化的人。

Most people would describe themselves as traders or investors, although there are no
consistently accepted definitions of either group. What I have to say is applicable to both
kinds of readers so I use the terms trading and investing interchangeably; the use of one
usually implies the other is included.

大多数人会把自己归类为交易者或投资者，尽管这两个群体并没有被普遍一致接受的严格定义。我在书中要讲的内容同时适用于这两类读者，
因此我会把交易和投资两个词交替使用；通常提到其中一个，也意味着另一个一并包含在内。

This book will be useful to amateurs --- individual investors trading their own money ---
and to market professionals who invest on behalf of others. The term ``amateur'' is not
intended to be patronising. It means only that you are not getting paid to manage money
and is no reflection on your level of skill.

本书既适合业余人士，也就是用自己的资金进行交易的个人投资者，也适合代表他人进行投资的市场专业人士。这里的“业余”一词并无任何居高临下之意，
它仅仅表示你并不是拿薪酬来管理资金，与能力高低无关。

This is not intended to be a parochial book solely for UK or US investors and I use
examples from a range of countries. Many books are written specifically for particular
asset markets. My aim is to provide a general framework that will suit traders of every
asset. I use specific examples from the equity, bond, foreign exchange and commodity
markets. These are traded with spread bets, exchange traded funds\footnote{All terms in
bold are defined in the glossary.\par 所有以黑体显示的术语都在词汇表中有定义。} and futures. But I do not explain the mechanics of
trading in detail. If you are not familiar with a particular market you should consult other
books or websites before designing your trading system.

这本书并不是只写给英国或美国投资者的狭义读物，我会使用来自多个国家的例子。许多书都只针对某一类特定资产市场而写，而我的目标是提供一个通用框架，
使其适用于各种资产的交易者。我会使用股票、债券、外汇和大宗商品市场中的具体例子，这些资产会通过点差投注、交易所交易基金和期货来交易。
不过，我不会详细讲解这些交易工具的具体操作机制。如果你对某个市场并不熟悉，那么在设计自己的交易系统之前，应该先去参考其他书籍或网站。

It might surprise you, but this book will also be useful for those who are sceptical of
computers entirely replacing human judgment. This is because there are several parts to a
complete trading system. Trading rules provide a prediction on whether something will go
up or down in price. These can be purely systematic, or based on human discretion. But
it is equally important to have a good framework of position and risk management. I
believe that a systematic framework should be used by all traders and investors for
position and risk management, even if the adoption of fully systematic rules is not
desirable.

你也许会感到意外，但这本书同样适用于那些对“让计算机完全取代人的判断”持怀疑态度的人。原因在于，一套完整的交易系统由多个部分组成。
交易规则负责预测某项资产的价格会上涨还是下跌，这些规则既可以完全系统化，也可以建立在人为裁量之上。但同样重要的是，要拥有一套完善的仓位与风险管理框架。
我相信，即使并不适合采用完全系统化的交易规则，所有交易者和投资者也都应该使用系统化框架来进行仓位与风险管理。

If you can beat simple rules when it comes to predicting prices, I show you how to use
your opinions in a systematic framework to make the best use of your talents.
Alternatively you might feel it is unlikely that anyone, man or machine, can predict the
markets. In this case, the same framework can be used to construct the best portfolio
consistent with that pessimistic view.

如果在价格预测这件事上，你确实能够胜过简单规则，那么我会告诉你，如何把你的判断纳入系统化框架中，以最大限度发挥你的优势。反过来说，
你也可能认为无论是人还是机器，都不太可能预测市场。在这种情况下，同样的框架也可以用来构建与这种悲观观点相一致的最佳投资组合。

\section*{Three examples}
\phantomsection
\addcontentsline{toc}{section}{Three examples}
\bisectiontitle{三个例子}

Throughout this book I focus on three typical groups of systematic traders and investors.
Don't panic if you don't fit neatly into any of these categories. I've chosen them because
between them they illustrate the most important issues which face all potential systematic
traders and investors.

在全书中，我会重点讨论三类典型的系统化交易者和投资者。如果你并不能完全归入其中任何一类，也不必紧张。我选择这三类人，是因为它们合起来，
能够展示所有潜在系统化交易者与投资者所共同面对的最重要问题。

In the final part of the book I discuss how to create a system tailored for each of these
audiences. Before then, each time you see one or more of these heading boxes it indicates
that the material in that section of the book is aimed mainly at the relevant group and is
optional for others.

在本书的最后一部分，我会讨论如何为这三类读者分别构建量身定制的系统。在此之前，每当你看到一个或多个此类标题框时，就表示该部分内容主要面向相应群体，
对其他读者而言则属于可选阅读。

\subsection*{Asset allocating investor}
\bisubsectiontitle{资产配置型投资者}

An asset allocating investor allocates funds amongst, and within, different asset classes.
Asset allocators can use systematic methods to avoid the short-term chasing of fads and
fashions that they know will reduce their returns. They might be lazy and wise amateur
investors, or managing institutional portfolios with long horizons such as pension funds.
Asset allocators are sceptical about those who claim to get extra returns from frequent
trading. For this reason the basic asset allocation example assumes you can't forecast how
asset prices will perform. However some investors might want to incorporate their views,
or the views of others. I show you how to achieve this without overtrading or ending up
with an extreme portfolio.

资产配置型投资者会在不同资产类别之间以及各类别内部进行资金分配。资产配置者可以借助系统化方法，避免追逐那些他们明知会损害收益的短期潮流与风尚。
他们可能是聪明且懂得偷懒的个人投资者，也可能是在管理像养老基金这类长期机构投资组合的专业人士。资产配置者通常对“频繁交易能带来额外收益”这类说法持怀疑态度。
因此，在基础资产配置的例子中，我会假设你无法预测资产价格的表现。不过，有些投资者可能希望纳入自己的观点，或者纳入他人的观点。我会展示如何在不过度交易、
也不把投资组合推向极端的前提下做到这一点。

Unlike the other examples asset allocators usually don't use leverage. I illustrate the
investment process with the use of unleveraged passive exchange traded funds (ETFs).
But the methods I show apply equally well to investors in collective funds of both the
active and passive varieties, and to those investing in portfolios of individual securities.
My own portfolio includes a basket of ETFs which I manage using the principles of the
asset allocating investor.

与其他例子不同，资产配置者通常不使用杠杆。我会用不加杠杆的被动型交易所交易基金（ETF）来说明这一投资流程。不过，我展示的方法同样适用于投资主动型或被动型集合基金的人，
以及投资于个别证券组合的人。我的个人投资组合中就包含一篮子 ETF，而我正是按照资产配置型投资者的原则来管理它们。

\subsection*{Semi-automatic trader}
\bisubsectiontitle{半自动交易者}

Semi-automatic traders live in a world of opportunistic bets\footnote{I am not using the
gambling term ``bet'' here in a pejorative sense. In my opinion the distinction some
people draw between financial gambling, trading and investing is completely meaningless:
they all involve taking financial risk on uncertain outcomes. Indeed professional gamblers
usually have a better understanding of risk management than many people working in the
investment industry.\par 我在这里使用“bet（下注）”这个词，并没有任何贬义。在我看来，有些人试图区分金融赌博、交易和投资，这种区分完全没有意义：
它们都意味着在不确定结果上承担金融风险。事实上，职业赌徒对风险管理的理解，往往比许多投资行业从业者更好。} taken on a fluid set of assets. Semi-automatic traders think they
are superior to simple rules when it comes to forecasting by how much prices will go up
or down; instead they make their own educated guesses. However they would like to
place those bets inside a systematic framework which will ensure their positions and risk
are properly managed. This frees them up to spend more time making the right call on the
market.

半自动交易者活跃在一个以机会性下注为主、且资产集合不断变化的世界里。半自动交易者认为，在判断价格会涨跌多少这件事上，自己比简单规则更高明，
因此他们会做出自己的理性判断。不过，他们仍希望把这些下注放进一个系统化框架中，以确保仓位与风险都得到妥善管理。这样一来，他们就能腾出更多时间，
专注于对市场方向作出正确判断。

In my example the semi-automatic trader is comfortable with leverage and investing with
derivatives. They are both buyers and sellers, betting for or against asset prices. My
semi-automatic trader is active in equity index and commodity spread bet markets, but
the example is widely applicable elsewhere.

在我的例子中，半自动交易者能够熟练使用杠杆和衍生品进行投资。他们既会做多，也会做空，也就是既可能押注资产价格上涨，也可能押注资产价格下跌。
我的半自动交易者活跃于股票指数和大宗商品点差投注市场，但这个例子的适用范围远不止于此。

I trade a portfolio of UK equities using the framework I've outlined here for
semi-automatic traders.

我自己就在用这里为半自动交易者勾勒出的这套框架，交易一个英国股票投资组合。

\subsection*{Staunch systems trader}
\bisubsectiontitle{坚定的系统化交易者}

The staunch systems trader is a true believer in the benefits of fully systematic trading.
Unlike the semi-automatic trader and the asset allocating investor, they embrace the use
of systematic trading rules to forecast price changes, but within the same common
framework for position risk management.

坚定的系统化交易者是真正相信完全系统化交易之益处的人。与半自动交易者和资产配置型投资者不同，他们愿意使用系统化交易规则来预测价格变化，
但仍会把这些规则置于同一个共同的仓位风险管理框架之中。

Many systems traders think they can find trading rules that give them extra profits, or
alpha. Others are unconvinced they have any special skill but believe there are additional
returns available which can't be captured just by ``buy and hold'' investing. They can use
very simple rules to capture these sources of alternative beta.

许多系统化交易者认为，他们可以找到能带来额外利润，也就是 alpha 的交易规则。另一些人则并不确信自己拥有什么特殊技能，但他们相信，市场中存在一些
仅靠“买入并持有”无法获得的额外收益，而这些收益可以通过非常简单的规则来捕捉，也就是另类贝塔的来源。

Systems traders may have access to back-testing software, either in off-the-shelf
packages, spreadsheets or bespoke software. It isn't absolutely necessary to have such
programs as I will be providing a flexible pre-configured trading system which doesn't
need back-testing. However if you want to develop your own new ideas I will show you
how to use these powerful software tools safely.

系统化交易者可能可以使用回测软件，无论是现成的软件包、电子表格，还是定制软件。拥有这些程序并非绝对必要，因为我会提供一套灵活的、预先配置好的交易系统，
它并不需要回测。不过，如果你想发展自己的新想法，我也会告诉你如何安全地使用这些强大的软件工具。

Like semi-automatic traders, staunch systems traders are comfortable with derivatives and
leverage. Although the examples I give are for futures trading, they are equally valid for
trading similar assets.

和半自动交易者一样，坚定的系统化交易者也能够熟练使用衍生品和杠杆。虽然我举的例子以期货交易为主，但这些原则同样适用于交易其他相似资产。

I trade over 40 futures contracts with my own money using a fully systematic set of eight
trading rules.

我用自己的资金交易 40 多个期货合约，并使用一套完全系统化、包含八条交易规则的方法来执行。

\section*{The technical stuff}
\phantomsection
\addcontentsline{toc}{section}{The technical stuff}
\bisectiontitle{技术性内容}

Inevitably this is a subject which requires some specialised terminology. Although I try
and keep jargon to a minimum it's usually easier to use a well-known shorthand term
rather than spelling everything out. Words and phrases highlighted in bold are briefly
defined in the glossary.

不可避免地，这是一个需要一些专业术语的主题。尽管我尽量把术语控制在最低限度，但通常使用大家熟悉的简称，会比把所有内容完整展开更方便。
文中以黑体标出的词语和短语，都会在词汇表中做简要定义。

As well as standard finance vocabulary I use my own invented terms. So phrases like
instrument block are also in bold type, and appear in the glossary with a short
explanation.

除了标准金融词汇之外，我还会使用自己创造的一些术语。因此，像 instrument block 这样的短语也会以黑体出现，并在词汇表中附有简短解释。

Certain important concepts require deeper understanding and I will include more detail
when I first use them, as in the box below.

某些重要概念需要更深入的理解，因此在它们第一次出现时，我会给出更详细的说明，就像下面这个框中展示的那样。

All detailed explanations in the text are signposted from the glossary, to help you refer to
them later in the book.

正文中的所有详细解释，都会在词汇表中给出索引提示，方便你在后文阅读时再回来查阅。

\conceptbox{%
\textbf{CONCEPT: EXAMPLE}

\medskip
These concept boxes give detailed explanations of key concepts in the book.

\medskip
\textbf{概念：示例}

\medskip
这些概念框会对本书中的关键概念作出更详细的说明。
}

\section*{What is coming}
\phantomsection
\addcontentsline{toc}{section}{What is coming}
\bisectiontitle{接下来会讲什么}

Running a systematic strategy is just like following any list of instructions, such as a
recipe. Many books on trading are like fast food outlets which give you something quick
and convenient to eat from a limited menu. However I am going to help you create your
own strategies; this is more than just a cookbook, it's a guide to writing recipes from
scratch. Inventing your own system requires extra work upfront, but it is more satisfying
and profitable in the long run.

执行一套系统化策略，就像遵循一份说明清单，比如菜谱。很多交易类书籍就像快餐店，只能从有限菜单里给你一份快速而方便的现成食物。
而我想帮助你做的是创建你自己的策略；这不只是一本菜谱，更是一份教你从零写出菜谱的指南。设计自己的系统，前期确实需要额外投入，
但从长远看，这会更令人满意，也更有盈利性。

To create your own recipes requires an understanding of the science of food chemistry and
of different kinds of dishes. Similarly, part one of this book --- ``Theory'' --- provides a
theoretical basis for why you should run systematic strategies and gives an overview of the
trading styles that are available.

要写出自己的菜谱，你需要理解食物化学原理以及不同菜式的特点。同样地，本书第一部分“理论”会提供一个理论基础，说明为什么你应该采用系统化策略，
并概览有哪些可行的交易风格。

Part two --- ``Toolbox'' --- provides you with two key methods used in the creation of
systematic strategies: back-testing and portfolio optimisation. Like sharp kitchen knives
these are powerful tools, but also potentially dangerous. When misused large trading
losses can be made despite apparently promising ideas. I will show you how to use them
properly --- and when you don't need them.

第二部分“工具箱”会提供两种构建系统化策略的关键方法：回测与投资组合优化。它们就像锋利的厨刀，是非常有力的工具，但也可能带来危险。
如果使用不当，即便想法看起来很有前景，也可能造成巨大的交易亏损。我会告诉你如何正确使用它们，以及在什么情况下其实并不需要它们。

Part three --- ``Framework'' --- provides a complete and extendable framework for the
creation of systematic strategies. Finally in part four --- ``Practice'' --- I show three
different uses of the framework in action. I illustrate how it can be adapted for each of
the semi-automatic trader, asset allocating investor and staunch systems trader.

第三部分“框架”会提供一个完整且可扩展的系统化策略构建框架。最后，在第四部分“实践”中，我会展示这个框架的三种不同应用方式，
说明它如何分别适配半自动交易者、资产配置型投资者和坚定的系统化交易者。

This book couldn't possibly be a comprehensive guide to the entire subject of trading.
Appendix A contains some books I would recommend for further reading and others I've
referred to in the text. There is also advice on where you can get data and access to
suitable brokers to begin investing systematically.

这本书不可能成为覆盖整个交易主题的百科全书。附录 A 列出了我推荐进一步阅读的一些书籍，以及正文中提到的其他资料。书中还会提供建议，
告诉你从哪里可以获取数据，以及如何接触到适合开始系统化投资的经纪商。

I've avoided putting mathematical formulas and detailed algorithms into the main text.
Instead appendix B contains detailed specifications for the trading rules, appendix C
covers portfolio optimisation, and appendix D includes further details on implementing
the framework. The website for this book --- \url{www.systematictrading.org} --- includes
additional material.

我尽量避免把数学公式和详细算法直接塞进正文。相应地，附录 B 包含交易规则的详细说明，附录 C 讨论投资组合优化，附录 D 则给出实施这一框架的更多细节。
本书网站 --- \url{www.systematictrading.org} --- 还提供额外材料。

This is not a book about automating trading strategies. It's possible to trade
systematically using an entirely manual process with just a spreadsheet to speed up
calculations, so automation is not necessary. Nevertheless automation is desirable when
running fast and complex strategies. My website also includes some details on my own
automated system and guidance to help you develop your own.

这不是一本关于“如何把交易策略自动化”的书。仅仅借助一个电子表格来加快计算，你也完全可以用纯手工流程进行系统化交易，因此自动化并非必要条件。
不过，当你运行速度快、结构复杂的策略时，自动化仍然是很有价值的。我的网站上也包含一些关于我自己自动化系统的细节，以及帮助你开发个人系统的指导。
